\begin{table}[!htbp]
\centering
\caption{Sensibilidade exploratória \textit{post hoc} da fusão: diferenças de revisões no \textit{bootstrap} e esforço dentro da faixa operacional. Valores negativos favorecem F4.}
\label{tab:cap8-fusao-robustez}
\small
\begin{tabular}{lrrr}
\toprule
Comparação & Meta & Mediana & IC 95\% \\
\midrule
F4 menos O1 & 9 & $-38$ & $-61$ a $-18$ \\
F4 menos O1 & 23 & $-274$ & $-995$ a $+40$ \\
F4 menos O1 & 36 & $-1.203$ & $-2.049$ a $-63$ \\
F4 menos O3 & 9 & $-17$ & $-44$ a $0$ \\
F4 menos O3 & 23 & $-25$ & $-402$ a $+59$ \\
F4 menos O3 & 36 & $-1.928$ & $-24.718$ a $+440$ \\
\midrule
\multicolumn{4}{l}{Revisões na faixa de 2.522 candidatos e 36 PAR~=~2} \\
Fila & Meta 9 & Meta 23 & Meta 36 \\
\midrule
O1 & 51 & 501 & 2.521 \\
O3 & 12 & 47 & 1.330 \\
F4 & 10 & 87 & 840 \\
F3 & 10 & 38 & 493 \\
\bottomrule
\end{tabular}
\par\footnotesize Fonte: \path{data/fusao_robustez_bootstrap.csv} e \path{data/fusao_robustez_faixa.csv}. Medianas e limites arredondados ao inteiro mais próximo. Reamostragem agrupada por óbito (2.000 réplicas), sem reajuste do classificador, com probabilidades fora da amostra e referência fixas. Metas de 9, 23 e 36 quando há pelo menos 36 PAR~=~2; nas demais réplicas, teto de 25\%, 64\% e 100\% do total, mudando o estimando. Filas e fusão recalculadas dentro da faixa.
\end{table}
